\section{MPC Attestation}


Messages are not self-authenticating. An MPC committee of $n$ nodes provides threshold attestation.

\subsection{Signing Protocol}

The MPC committee uses CGGMP21~\cite{cggmp2021} for ECDSA threshold signing (EVM-compatible chains) and FROST~\cite{frost2020} for Schnorr threshold signing (non-EVM chains).

For a message $M$ with hash $H = \text{keccak256}(\text{abi.encode}(M))$:

\begin{enumerate}[leftmargin=2em]
\item Each MPC node independently verifies the source chain event (deposit confirmation, finality).
\item Each node that confirms the event produces a signature share over $H$.
\item The coordinator collects $t$-of-$n$ shares and produces the threshold signature $\sigma$.
\item The message $(M, \sigma)$ is submitted to the destination chain.
\end{enumerate}

\subsection{Verification on Destination}

The destination chain verifier holds the MPC committee's public key (registered on the P-Chain). Verification is a single \texttt{ecrecover} (ECDSA) or Schnorr verification---identical to verifying a single signature. The threshold structure is invisible to the verifier.

\subsection{Key Rotation}

MPC key rotation occurs via DKG on the M-Chain:
\begin{enumerate}[leftmargin=2em]
\item New key shares are generated via proactive secret sharing.
\item The new public key is registered on the P-Chain.
\item A transition period accepts signatures from both old and new keys.
\item After the transition period, old keys are deprecated.
\end{enumerate}

